% ===================================================================
% MT-04b: Minitest Stundenplan (El horario) - MUSTERLÖSUNG
% Spanisch Klasse 7
% ===================================================================

\documentclass[11pt, a4paper]{article}

\usepackage[utf8]{inputenc}
\usepackage[T1]{fontenc}
\usepackage[ngerman]{babel}
\usepackage{helvet}
\renewcommand{\familydefault}{\sfdefault}

\usepackage[margin=2cm]{geometry}
\usepackage{enumitem}
\usepackage{tabularx}
\usepackage{booktabs}

\usepackage{xcolor}
\usepackage{tcolorbox}
\tcbuselibrary{skins}

\usepackage{fancyhdr}
\usepackage{lastpage}

% FARBEN
\definecolor{spanisch}{HTML}{c41e3a}
\definecolor{grau}{HTML}{888888}
\definecolor{hellgrau}{HTML}{f5f5f5}
\definecolor{loesung}{HTML}{c62828}

\newcommand{\fachfarbe}{spanisch}
\newcommand{\thema}{Minitest: El horario -- MUSTERLÖSUNG}
\newcommand{\fach}{Spanisch}
\newcommand{\klasse}{7}
\newcommand{\maxpunkte}{28}
\newcommand{\zeit}{15 Minuten}
\newcommand{\datum}{\today}

\pagestyle{fancy}
\fancyhf{}
\renewcommand{\headrulewidth}{0.4pt}
\renewcommand{\footrulewidth}{0.4pt}

\lhead{\fach{} -- Klasse \klasse}
\chead{\textbf{MT-04b -- ML}}
\rhead{\datum}

\lfoot{\zeit}
\cfoot{}
\rfoot{Seite \thepage\ von \pageref{LastPage}}

\newcommand{\punktebox}[1]{\hfill\fbox{\small \_\_\_/#1}}
\newcommand{\luecke}[1]{\rule{#1}{0.4pt}}
\newcommand{\lsg}[1]{\textcolor{loesung}{\textbf{#1}}}

\newtcolorbox{infobox}{
    colback=hellgrau,
    colframe=\fachfarbe,
    boxrule=1pt,
    arc=0pt,
    left=8pt, right=8pt, top=8pt, bottom=8pt
}

\newenvironment{aufgabe}[1]{%
    \vspace{12pt}
    \noindent\textbf{\textcolor{\fachfarbe}{Aufgabe #1}}
    \vspace{6pt}
    \begin{list}{}{\leftmargin=0pt}
    \item[]
}{%
    \end{list}
}

\newcommand{\es}[1]{\textcolor{\fachfarbe}{\textbf{#1}}}

\begin{document}

% ===================================================================
% KOPFBEREICH
% ===================================================================

\begin{center}
    {\Large\bfseries\textcolor{loesung}{\thema}}
\end{center}

\vspace{5pt}

\noindent
\begin{tabularx}{\textwidth}{|X|X|X|}
\hline
\textbf{Name:} \lsg{LÖSUNG} &
\textbf{Klasse:} \klasse &
\textbf{Datum:} \luecke{2cm} \\
\hline
\end{tabularx}

\vspace{10pt}

\begin{infobox}
\textbf{Zeit:} \zeit{} | \textbf{Punkte:} \maxpunkte{} BE | \textbf{Hilfsmittel:} keine
\end{infobox}

\vspace{15pt}

% ===================================================================
% AUFGABE 1: Wochentage
% ===================================================================

\begin{aufgabe}{1 -- Wochentage in Reihenfolge}
Schreibe die Wochentage in der richtigen Reihenfolge auf Spanisch. \punktebox{7}
\end{aufgabe}

\noindent
\begin{tabular}{rlrl}
1. (Montag) & \lsg{lunes} & 2. (Dienstag) & \lsg{martes} \\[0.8em]
3. (Mittwoch) & \lsg{miércoles} & 4. (Donnerstag) & \lsg{jueves} \\[0.8em]
5. (Freitag) & \lsg{viernes} & 6. (Samstag) & \lsg{sábado} \\[0.8em]
7. (Sonntag) & \lsg{domingo} & & \\
\end{tabular}

\textcolor{loesung}{\textit{(7 BE: je 1 BE pro Wochentag)}}

\vspace{15pt}

% ===================================================================
% AUFGABE 2: Uhrzeiten
% ===================================================================

\begin{aufgabe}{2 -- Uhrzeiten auf Spanisch}
Schreibe die Uhrzeiten auf Spanisch. \punktebox{7}
\end{aufgabe}

\noindent
a) 8:00 Uhr: \lsg{Son las ocho.} \\[0.8em]
b) 10:30 Uhr: \lsg{Son las diez y media.} \\[0.8em]
c) 1:00 Uhr: \lsg{Es la una.} \\[0.8em]
d) 3:30 Uhr: \lsg{Son las tres y media.} \\[0.8em]
e) 12:00 Uhr: \lsg{Son las doce.} \\[0.8em]
f) 9:30 Uhr: \lsg{Son las nueve y media.} \\[0.8em]
g) 2:00 Uhr: \lsg{Son las dos.} \\

\textcolor{loesung}{\textit{(7 BE: je 1 BE pro Uhrzeit)}}

\textcolor{loesung}{\textit{Hinweis: Bei 1 Uhr = ``Es la una'' (Singular!). Sonst = ``Son las...'' (Plural). ``y media'' = +30 Minuten.}}

\vspace{15pt}

% ===================================================================
% AUFGABE 3: Fragen und Antworten
% ===================================================================

\begin{aufgabe}{3 -- Fragen und Antworten}
Ordne die passende Antwort der Frage zu. (A, B oder C) \punktebox{8}
\end{aufgabe}

\noindent
\begin{tabular}{|p{6cm}|p{6cm}|}
\hline
\textbf{Fragen} & \textbf{Antworten} \\
\hline
1. \es{¿Qué día es hoy?} \lsg{C} & A) Los viernes a las diez. \\[0.5em]
2. \es{¿A qué hora tienes español?} \lsg{B} & B) A las nueve. \\[0.5em]
3. \es{¿Cuándo tienes música?} \lsg{A} & C) Hoy es lunes. \\[0.5em]
4. \es{¿Qué hora es?} \lsg{D} & D) Son las once. \\
\hline
\end{tabular}

\textcolor{loesung}{\textit{(8 BE: je 2 BE pro richtige Zuordnung)}}

\textcolor{loesung}{\textit{Tipps: ¿Qué día? = Tag, ¿A qué hora? = Uhrzeit, ¿Cuándo? = Tag + Uhrzeit, ¿Qué hora es? = Uhrzeit}}

\vspace{15pt}

% ===================================================================
% AUFGABE 4: Sätze schreiben
% ===================================================================

\begin{aufgabe}{4 -- Schreibe auf Spanisch}
Übersetze die Sätze ins Spanische. \punktebox{6}
\end{aufgabe}

\noindent
a) Heute ist Montag. \\[0.5em]
\lsg{Hoy es lunes.} \\[1em]

\noindent
b) Es ist 10 Uhr. \\[0.5em]
\lsg{Son las diez.} \\[1em]

\noindent
c) Um 8 Uhr habe ich Mathe. \\[0.5em]
\lsg{A las ocho tengo matemáticas.} (oder: \lsg{Tengo matemáticas a las ocho.}) \\

\textcolor{loesung}{\textit{(6 BE: je 2 BE pro korrekten Satz)}}

% ===================================================================
% PUNKTEÜBERSICHT
% ===================================================================

\vspace{20pt}
\hrule
\vspace{10pt}

\noindent
\begin{tabularx}{\textwidth}{|X|c|c|c|c||c|}
\hline
\textbf{Aufgabe} & 1 & 2 & 3 & 4 & \textbf{Gesamt} \\
\hline
\textbf{max. Punkte} & 7 & 7 & 8 & 6 & \textbf{28} \\
\hline
\textbf{erreicht} & & & & & \\[0.8em]
\hline
\end{tabularx}

\vspace{15pt}

\noindent\textbf{Bewertungsschlüssel (14-NP):}

\begin{center}
\begin{tabular}{|c|c|c|c|c|c|c|c|c|c|c|c|c|c|c|}
\hline
\textbf{NP} & 14 & 13 & 12 & 11 & 10 & 9 & 8 & 7 & 6 & 5 & 4 & 3 & 2 & 1 \\
\hline
\textbf{\%} & 100 & 96 & 91 & 86 & 80 & 73 & 67 & 60 & 53 & 47 & 40 & 33 & 27 & 20 \\
\hline
\textbf{BE} & 28 & 27 & 25 & 24 & 22 & 20 & 19 & 17 & 15 & 13 & 11 & 9 & 8 & 6 \\
\hline
\end{tabular}
\end{center}

\end{document}
